\documentclass[a4paper,11pt]{bxjsarticle}
\usepackage[margin=25mm]{geometry}
\usepackage{amsmath,amssymb,booktabs}
\usepackage{hyperref}
\usepackage{xcolor}
\hypersetup{colorlinks=true,linkcolor=blue,urlcolor=blue}

\title{EpisodicDT-small: 学習手法とアルゴリズム}
\author{}
\date{\today}

\begin{document}
\maketitle

\section{概要}
本資料は、RSRP（Reference Signal Received Power）時系列から将来のRSRP系列を生成する
EpisodicDT-smallの学習・生成アルゴリズムを説明する。実装は、過去系列を
Transformer Encoderで文脈ベクトルへ圧縮し、その文脈を条件とする拡散モデルで将来系列を
生成する条件付き拡散モデルである。名称にDecision Transformerを含むが、現在の実装は
行動や報酬を入力とする標準的なDecision Transformerではなく、RSRPの確率的将来予測を
目的とするモデルである。

本資料中の行番号は、2026年9月7日時点のソースコードに対応する。

\section{データとエピソード}
\subsection{入力データと時系列分割}
4つのSUTD 5GシナリオCSVから\texttt{Time}と\texttt{RSRP}を抽出する。
各シナリオは時刻順に並べ、前半を学習用、後半を最終評価用にする。
前処理は\texttt{src/data/prepare\_sutd\_5g.py}に実装されている。
データは約2 Hzであり、1ステップは約0.5秒である。

各時系列$\{r_1,\ldots,r_N\}$から、過去長$H=20$、未来長$F=20$のエピソードを作る。
開始位置$s$のエピソードは
\begin{align}
  \mathbf{x}^{(s)} &= (r_s,\ldots,r_{s+H-1}), \\
  \mathbf{y}^{(s)} &= (r_{s+H},\ldots,r_{s+H+F-1})
\end{align}
である。したがって、historyとfutureは各10秒、1エピソードは約20秒に相当する。

RSRPWindowDatasetは各CSV内でのみウィンドウを作るため、異なるシナリオをまたぐ
エピソードは作られない。対応箇所は\texttt{src/data/dataset.py}の34--38行目である。
エピソードのhistoryとfutureの切り出しは43--50行目で行われる。

\subsection{標準化}
学習用トレース全体の平均$\mu$、標準偏差$\sigma$を用いてRSRPを標準化する。
\begin{equation}
  \tilde{r} = \frac{r-\mu}{\sigma}.
\end{equation}
標準化統計量の計算は\texttt{dataset.py}の30--33行目、各エピソードへの適用は48--49行目である。
評価・出力時には
\begin{equation}
  r = \tilde{r}\sigma + \mu
\end{equation}
によりdBm単位へ戻す。

\subsection{学習・検証・評価の分離}
最終評価用の後半データは学習に使わない。さらに、学習用前半データの各シナリオについて、
時系列後半20\%を検証用にする。検証境界をまたぐ40ステップのエピソードは除外するため、
学習・検証間で同一観測点を共有するエピソードはない。

この時系列分割は\texttt{src/train/train.py}の16--33行目で実装されている。
設定値\texttt{validation\_fraction}を使う呼出しは65--70行目である。

\section{モデル}
\subsection{Transformer Encoderによるエピソード表現}
historyの標準化系列を
\begin{equation}
  \tilde{\mathbf{x}} \in \mathbb{R}^{B\times H\times 1}
\end{equation}
とする。$B$はバッチサイズである。まず入力値を$D=128$次元へ線形変換し、位置埋め込みを加える。
\begin{equation}
  \mathbf{h}^{(0)}_i = W_{\mathrm{in}}\tilde{x}_i + \mathbf{p}_i.
\end{equation}
位置埋め込みはsin/cos形式であり、次式で表される。
\begin{align}
  p_{i,2k} &= \sin\left(i / 10000^{2k/D}\right), \\
  p_{i,2k+1} &= \cos\left(i / 10000^{2k/D}\right).
\end{align}

次に$L$層のTransformer Encoderを適用する。
\begin{equation}
  \mathbf{H}^{(L)} = \mathrm{TransformerEncoder}(\mathbf{H}^{(0)}).
\end{equation}
最終時刻のトークンを32次元へ射影して、future生成の条件$c$とする。
\begin{equation}
  \mathbf{c}=\mathrm{SiLU}(W_c\mathbf{h}^{(L)}_{H}+b_c),
  \qquad \mathbf{c}\in\mathbb{R}^{32}.
\end{equation}

実装の対応は表\ref{tab:encoder}の通りである。
\begin{table}[htbp]
\centering
\caption{Transformer Encoderの実装対応}
\label{tab:encoder}
\begin{tabular}{ll}
\toprule
コード箇所 & 処理 \\
\midrule
\texttt{episodic\_diffusion.py:12--22} & 入力射影、Transformer層、context射影の定義 \\
\texttt{episodic\_diffusion.py:27--32} & sin/cos位置埋め込みの計算 \\
\texttt{episodic\_diffusion.py:34--37} & Transformer適用と最終トークンのcontext化 \\
\bottomrule
\end{tabular}
\end{table}

現在の設定では$D=128$、attention head数は4、Transformer層数は3、dropoutは0.05である。

\subsection{条件付きノイズ予測器}
ノイズ予測器は、ノイズ付きfuture $\mathbf{y}_t$、正規化された拡散時刻$t$、
および文脈$c$を受け取るMLPである。
\begin{equation}
  \hat{\boldsymbol{\epsilon}}=
  \epsilon_\theta(\mathbf{y}_t,t,\mathbf{c}).
\end{equation}
実装では、futureを長さ20のベクトルに平坦化し、contextと時刻埋め込みを連結する。
\begin{equation}
  \hat{\boldsymbol{\epsilon}}=
  \mathrm{MLP}([\mathrm{flatten}(\mathbf{y}_t);\mathbf{c};\mathrm{TimeEmbed}(t)]).
\end{equation}
定義は\texttt{episodic\_diffusion.py}の40--60行目である。

\section{拡散学習}
\subsection{順方向拡散}
futureの正解系列を$\tilde{\mathbf{y}}$とする。各ミニバッチの各サンプルに対して、
拡散時刻$t\in\{0,\ldots,T-1\}$を一様に選ぶ。現在の設定は$T=100$である。

線形スケジュール$\beta_t$を用い、
\begin{align}
  \alpha_t &= 1-\beta_t, \\
  \bar{\alpha}_t &= \prod_{s=0}^{t}\alpha_s
\end{align}
と定義する。正規分布からノイズをサンプルする。
\begin{equation}
  \boldsymbol{\epsilon}\sim\mathcal{N}(\mathbf{0},I).
\end{equation}
ノイズ付きfutureは
\begin{equation}
  \mathbf{y}_t=
  \sqrt{\bar{\alpha}_t}\tilde{\mathbf{y}}+
  \sqrt{1-\bar{\alpha}_t}\boldsymbol{\epsilon}
\end{equation}
で得られる。

\texttt{episodic\_diffusion.py}の74--78行目がスケジュールを、80--89行目が時刻・
ノイズ・ノイズ付きfutureの作成および予測器の呼出しを実装している。

\subsection{損失関数と更新}
損失は、実際に加えたノイズと予測ノイズの平均二乗誤差である。
\begin{equation}
  \mathcal{L}_{\mathrm{diffusion}}(\theta)=
  \frac{1}{BF}\sum_{b=1}^{B}\sum_{j=1}^{F}
  \left(\epsilon_{b,j}-\hat{\epsilon}_{\theta,b,j}\right)^2.
\end{equation}
これは\texttt{episodic\_diffusion.py:90}の\texttt{mse\_loss}に対応する。
エンコーダ単独の損失はなく、この損失の逆伝播によりTransformer Encoderとノイズ予測器の
両方が学習される。

最適化にはAdamWを使用する。学習ループの対応箇所は表\ref{tab:training}の通りである。
\begin{table}[htbp]
\centering
\caption{学習ループの実装対応}
\label{tab:training}
\begin{tabular}{ll}
\toprule
コード箇所 & 処理 \\
\midrule
\texttt{train.py:76} & AdamWの生成（学習率$10^{-4}$） \\
\texttt{train.py:85--86} & ミニバッチの取得とデバイス転送 \\
\texttt{train.py:87--88} & 勾配初期化と拡散損失の計算 \\
\texttt{train.py:89} & 誤差逆伝播 \\
\texttt{train.py:90} & 勾配ノルムを1.0にクリッピング \\
\texttt{train.py:91} & AdamWによるパラメータ更新 \\
\bottomrule
\end{tabular}
\end{table}

\section{Early Stoppingと最良モデルの選択}
各エポックの学習後、検証データに対する平均拡散MSEを計算する
（\texttt{train.py:36--43, 93}）。検証MSEが
\texttt{early\_stopping\_min\_delta}=$10^{-4}$以上改善したとき、そのエポックの重みを最良として保存する
（98--102行目）。改善しないエポック数を数え（103--104行目）、最低30エポックを過ぎ、
15エポック連続で改善しなければ学習を終了する（105--107行目）。

最後に最良の重みを復元して保存する（109--120行目）。そのため、\texttt{epochs}=150は
最大エポック数であり、実際の終了エポックは検証損失により自動決定される。

\section{生成（逆拡散）}
学習後、historyからcontextを一度だけ計算する。次に、futureと同じ形状のガウスノイズを
初期状態$\mathbf{x}_T$とする。
\begin{equation}
  \mathbf{x}_T\sim\mathcal{N}(\mathbf{0},I).
\end{equation}
各$t=T-1,\ldots,0$に対してノイズを予測し、次式で逆拡散する。
\begin{equation}
  \mathbf{x}_{t-1}=
  \frac{1}{\sqrt{\alpha_t}}
  \left(
    \mathbf{x}_{t}-
    \frac{1-\alpha_t}{\sqrt{1-\bar{\alpha}_t}}
    \hat{\boldsymbol{\epsilon}}_t
  \right)
  + \mathbb{I}[t>0]\sqrt{\beta_t}\mathbf{z},
  \qquad \mathbf{z}\sim\mathcal{N}(\mathbf{0},I).
\end{equation}
この処理は\texttt{episodic\_diffusion.py:92--111}に実装されている。
初期ノイズを変えることで、同じhistoryに対して複数の将来RSRP候補を生成できる。

\section{最終評価と可視化}
\texttt{evaluate.py:16--26}では、最終評価データのすべてのエピソードに対して複数のfutureを
生成し、標準化を解除する。全エピソードの生成値はCSVに保存される
（\texttt{evaluate.py:136--154}）。

\texttt{forecast.png}は、評価エピソードから均等に選んだ代表例をsmall multiplesとして表示する。
各パネルは、history、実測future、生成中央値、生成分布の50\%区間と90\%区間、危険閾値を含む。
描画処理は\texttt{evaluate.py:54--91}である。さらに、全評価エピソード・全future stepを
集約した実測分布と生成分布の箱ひげ図を\texttt{forecast\_boxplot.png}に保存する
（\texttt{evaluate.py:29--51}）。

\end{document}
